\section{Discussione}
\label{sec:discussion}

La discussione sarà organizzata intorno alla distinzione fra equivalenza
semantica e prestazione dell'esecutore. Per le modalità sincrone, ogni
accelerazione dovrà essere interpretata alla luce del controllo di traiettoria
esatta e del costo della funzione obiettivo. Per asincronia e isole, le
variazioni della qualità saranno discusse come effetti di una diversa dinamica
di aggiornamento e comunicazione, non come accelerazioni del medesimo
algoritmo.

\subsection{Implicazioni per la parallelizzazione CPSO}

Da completare con il rapporto fra granularità del lavoro, overhead di
scheduling, sincronizzazione, dimensione della popolazione e costo della
valutazione. Il confronto con PAPSO, le implementazioni CPU/GPU e i recenti
lavori multicore dovrà essere qualitativo e contestualizzato, senza confrontare
tempi assoluti ottenuti su hardware, algoritmi o protocolli differenti.

\subsection{Ruolo delle topologie e della migrazione}

Da completare distinguendo qualità, densità/connettività, grado medio, costo
della selezione dei vicini e frequenza delle comunicazioni. In assenza di un
disegno degree-matched, non attribuire causalmente le differenze alla sola
geometria del grafo.

\subsection{Confronto con le baseline e validità esterna}

Da completare discutendo gli esiti per famiglia di benchmark e meccanismo PSO.
Il confronto rimane limitato alle dieci baseline PSO implementate e non
permette conclusioni di superiorità rispetto a CMA-ES, differential evolution
o all'intero stato dell'arte dell'ottimizzazione continua.

\subsection{Minacce alla validità e limiti}

Da documentare: numero di funzioni e unità statistiche; hardware locale
fanless e variabilità termica; differenze di generatore casuale e
inizializzazione fra implementazioni; adattamenti delle baseline per FE parity;
callback Python seriali nei benchmark esterni; differenze fra pilot locale e
campagne HPC; generalizzazione a obiettivi costosi, GPU e sistemi distribuiti.
